% README_charge_distribution.tex
% Description: LaTeX-formatted explanation of the charge-distribution
% algorithm used in the DustRAMSES repository and the memory
% optimizations implemented. Includes pseudocode suitable for
% inclusion in a paper.

\documentclass{article}
\usepackage{amsmath,amssymb}
\usepackage{booktabs}
\usepackage{listings}
\usepackage{graphicx}
\usepackage{hyperref}
\usepackage{caption}
\usepackage{geometry}
\geometry{margin=1in}

\begin{document}

	itle{Computation of Equilibrium Grain Charge Distributions\and Memory Optimizations (DustRAMSES)}
\author{}
\date{}
\maketitle

\section*{Abstract}
This note provides a concise, reproducible description of the numerical
methods implemented in \texttt{DustRAMSES} to compute equilibrium grain charge
distributions. It documents the underlying physical model, the numerical
algorithms used to obtain stable and accurate probability distributions
P(Z), and the memory- and performance-related optimizations employed to make
large-scale evaluations tractable. The exposition is written to be suitable
for inclusion in a methods section of a manuscript: it contains the key
equations, implementation-level pseudocode, and a discussion of edge cases,
complexity, and recommended operating regimes.

\section{Physical model and rates}
We compute the equilibrium probability distribution $P(Z)$ for a grain to
possess integer charge $Z$ by enforcing detailed balance among the
microphysical charging channels. The dominant processes included in the
present implementation are:
\begin{itemize}
  \item Photoelectric emission (rate per grain):
  \begin{equation*}
    R_{\rm pe}(Z) = \int J_{\nu}(\nu)\; Y(Z,E)\; C_{\rm abs}(\nu,a)\; F(E,Z)\,d\nu.
  \end{equation*}
  Here $J_{\nu}$ is the specific radiation intensity (photons or energy as
  appropriate for the chosen units), $C_{\rm abs}(\nu,a)$ is the grain
  absorption cross section for radius $a$, $Y(Z,E)$ is the photoelectric yield
  as a function of emitted-electron energy $E$ and grain charge $Z$, and
  $F(E,Z)$ denotes the fraction of the absorbed photon energy available to
  liberate an electron (this factor depends on the material work function and
  the grain electrostatic potential). The treatment of photoelectric yields
  and the associated kinetic-energy distributions follows the formalism of
  Weingartner & Draine (2001) and references therein \cite{Weingartner2001}.
  \item Electron recombination and collisional charging by ions and electrons:
  the rates $R_{\rm col}^{e,i}(Z)$ are computed from collisional cross
  sections folded with the Maxwellian velocity distribution of the plasma
  species. In the low-energy (thermal) limit these rates scale proportionally
  to $n_{e}\,\langle v\rangle\,\sigma$, with appropriate Coulomb-focusing or
  suppression factors applied for charged grains.
  \item Secondary electron emission (including Auger-like processes) and other
  minor channels are included as additive contributions when the material
  model or energy regime warrants their inclusion.
\end{itemize}

For compact notation we denote the net upward transition rate (loss of one
electron) from $Z$ to $Z+1$ by $R^{+}(Z)$ and the downward transition rate
(gain of one electron) by $R^{-}(Z)$. At steady state the detailed-balance
recurrence yields
\begin{equation}
  P(Z+1) = P(Z)\,\frac{R^{+}(Z)}{R^{-}(Z+1)}
\end{equation}

Normalization is imposed by requiring $\sum_Z P(Z)=1$ after truncating the
charge domain to a finite bracket $Z_{\min}\ldots Z_{\max}$ chosen so that
the tail probabilities at the boundaries are below a user-specified tolerance
(e.g., $10^{-8}$). This truncation is standard practice for numerically
stable representations of discrete probability distributions.

\section{Numerical algorithm}
The numerical solver implemented in the code proceeds according to the
following high-level algorithmic steps:
\begin{enumerate}
  \item Precompute or interpolate the photon absorption cross section
    $C_{\rm abs}(\nu,a)$ for the grain material and radius of interest. The
    implementation supports interpolation from tabulated optical constants or
    direct evaluation from analytic dielectric models where available.
  \item Compute (or retrieve from cache) the photoelectric integrand
    $J_{\nu}\,Y(Z,E)\,C_{\rm abs}(\nu,a)$ over the radiation bins. For
    computational efficiency the code constructs this kernel in vectorized
    blocks over photon-energy bins and integrates (reduces) the contribution
    to $R_{\rm pe}(Z)$ for each $Z$. Where possible, kernels that are
    identical across multiple evaluation points (for example, identical
    radiation fields and cross sections) are reused to avoid redundant work
    (cache keyed by grain, radius, and radiation model).
  \item Compute collisional rates $R^{-}(Z)$ (electron and ion recombination)
    using analytic formulae augmented by interpolated stopping-power tables
    where necessary. Coulomb focusing and shielding corrections are applied as
    appropriate for the plasma conditions.
  \item Select a reference charge $Z_{\rm ref}$ close to the estimated mean
    charge, and determine a safe computational bracket $[Z_{\min},Z_{\max}]$
    by expanding until the tail probabilities fall below a prescribed
    tolerance (see below). This limits the number of discrete states that
    must be evaluated without compromising numerical accuracy.
  \item Use the detailed-balance recurrence to compute $P(Z)$ relative to a
    reference value $P(Z_{\rm ref})$. The recurrence is implemented in
    logarithmic form to preserve numerical stability when probabilities span
    many orders of magnitude (see Sec.~\ref{sec:stability}).
  \item Normalize the distribution and compute summary statistics: mean
    $\langle Z\rangle$, variance, and higher moments if desired.
\end{enumerate}

\subsection*{Choosing $Z_{\rm ref}$ and bounds}
The code uses the following pragmatic strategy to select $Z_{\rm ref}$ and the
computational bounds:
\begin{itemize}
  \item Estimate a reference mean charge from a fast, low-precision model
    (for instance by equating characteristic photoemission and collisional
    rates or by evaluating rates at representative environmental parameters).
  \item Set $Z_{\rm ref}=\mathrm{round}(Z_{\rm est})$ and expand outward
    until the cumulative tail mass on both sides is less than the target
    tolerance (e.g., $10^{-8}$), thus defining $Z_{\min}$ and $Z_{\max}$.
\end{itemize}

\section{Stability and log-recursion}
Rather than compute $P(Z)$ directly (which may underflow when probabilities
are extremely small) the code computes logarithmic probabilities using the
recurrence
\begin{equation*}
  \log P(Z+1) = \log P(Z) + \log R^{+}(Z) - \log R^{-}(Z+1).
\end{equation*}
Working in log space keeps intermediate values within the numeric range of
IEEE floating-point arithmetic and allows stable normalization via the
log-sum-exp operation. This approach follows best-practice numerical
techniques for discrete-state master equations and is used in related grain
charging literature (see, e.g., \cite{Weingartner2001}).

\section{Memory and performance optimizations}
The implementation targets broad parameter sweeps over grain sizes and
radiation fields, where per-combination cost can be high. The principal
memory- and performance-reduction strategies employed are:

\begin{enumerate}
  \item \textbf{Vectorized integrals with chunking:} constructing the full
    photoelectric kernel over an extensive photon-energy grid can produce very
    large temporary arrays. We therefore build the kernel in memory-limited
    blocks (chunks) over contiguous photon-energy bins, integrate each chunk
    and accumulate the partial contribution to $R_{\rm pe}(Z)$. This reduces
    peak memory usage while maintaining vectorized arithmetic and high
    throughput.

  \item \textbf{Per-worker temporary budget and streaming fallback:}
    For parallel execution the code estimates available memory and sets a
    conservative per-worker temporary budget. If the estimated memory
    footprint of a fully vectorized kernel exceeds the budget the worker
    automatically switches to a streaming (per-$Z$ or per-energy) integration
    loop that uses a small, constant memory footprint. The streaming mode is
    slower but avoids swapping and out-of-memory failures on memory-constrained
    systems.

  \item \textbf{Caching of photoelectric integrals:} computation of
    $R_{\rm pe}(Z)$ dominates runtime. The implementation therefore caches
    previously computed photoelectric integrals keyed by the combination of
    (material model, radius, radiation field, $Z$, yield parameters). The
    cache uses an LRU eviction policy and is intentionally shallow to avoid
    retaining very large arrays in memory for long periods.

  \item \textbf{Adaptive windowing for $\gamma$–$T$ sampling:} when
    mapping the dimensionless charging parameter $\gamma = G_0\sqrt{T}/n_e$ to
    computed charge distributions, the code groups results into logarithmic
    bins. The number of bins is set adaptively based on the number of unique
    $\gamma$ samples to reduce output size and to avoid storing redundant
    per-sample distributions when many combinations correspond to the same
    effective charging parameter.

  \item \textbf{Controlled parallelism:} parallel worker counts are chosen to
    respect the per-worker temporary budget and the host memory capacity, in
    order to avoid over-subscription and contention on shared systems.

  \item \textbf{Precision and dtype selection:} intermediate temporary
    buffers that are not precision-critical (for example, chunked integrand
    arrays) may be held in \texttt{float32} to reduce memory usage. Final
    reductions, accumulations of rates, and the log-probability representation
    are performed in \texttt{float64} to maintain numerical fidelity.

\end{enumerate}

\section{Pseudocode}
The following condensed pseudocode presents the essential implementation
steps used to compute the equilibrium distribution for a single tuple
$(G_0,T,n_e)$ or for a vectorized batch. It is intentionally implementation
focused to aid reproducibility and direct reuse in manuscript methods.

\begin{lstlisting}[language=Python, caption={Pseudocode for equilibrium charge distribution}]
# Inputs: grain_type, a_micron, radiation_field (J_nu, nu), T, ne, ion_species
# yield_params controls internal numerical budgets and caching

# 1) Prepare cross sections and yield kernels (interpolate if needed)
C_abs = load_or_interpolate_cross_section(grain_type, a_micron)
# If cached R_pe kernels exist for the same (grain,a,rad_model), reuse them

# 2) Compute collisional rates R_minus(Z) (electrons + ions)
R_minus = compute_collisional_rates(T, ne, ion_species, grain_type, a_micron)

# 3) Compute photoelectric upward rates R_plus(Z)
#    - build J_nu * C_abs * yield(E,Z) as integrand and integrate over nu
R_plus = zeros_like(Z_range)
for energy_chunk in energy_grid.chunks(max_bytes=worker_budget):
  integrand = J_nu[chunk] * C_abs[chunk] * yield_function(Z_range, energy_chunk)
  # integrate per Z across the chunk and accumulate
  R_plus += trapz(integrand, nu[chunk])

# 4) Choose Z_ref and bracket [Zmin,Zmax]
Z_ref = estimate_reference_charge(R_plus, R_minus)
Zmin, Zmax = expand_bounds_until_tail_small(Z_ref, R_plus, R_minus, tol=1e-8)
Z_range = arange(Zmin, Zmax+1)

# 5) Compute log-probabilities by recurrence from Z_ref outwards
logP = full_like(Z_range, -inf, dtype=float)
logP[Z_ref_index] = 0.0  # reference
# upward recursion
for i from Z_ref_index to end-1:
  logP[i+1] = logP[i] + log(R_plus[Z_i]) - log(R_minus[Z_i+1])
# downward recursion
for i from Z_ref_index down to 1:
  logP[i-1] = logP[i] + log(R_minus[Z_i]) - log(R_plus[Z_i-1])

# 6) Normalize using log-sum-exp
log_sum = logsumexp(logP)
P = exp(logP - log_sum)

# Output: P(Z), mean, sigma
Zmean = sum(Z_range * P)
Zvar = sum((Z_range - Zmean)**2 * P)

done


# Memory-safe variant: if vectorized kernels exceed budget, compute R_plus per Z
if kernel_bytes > worker_budget:
  for Z in Z_range:
    R_plus[Z] = integrate_over_nu(J_nu * C_abs * yield(Z, E))

done


# Batch mode: when evaluating many (G0,T,ne) combinations that share
# the same cross sections and radiation model, reuse C_abs and cached kernels
# and compute only deltas in the collisional rates.

done

\end{lstlisting}

\section{Edge cases and robustness}
\begin{itemize}
  \item \textbf{Single-charge outcomes:} when the distribution is essentially
    monodisperse (e.g., $P(Z_0)\approx 1$), the solver returns a delta-like
    distribution. The code may emit an informational message in this regime to
    flag the degenerate outcome.
  \item \textbf{Non-finite input values:} input tuples that produce
    non-finite rates (NaN or Inf) are removed prior to binning and aggregation.
    If no valid combinations remain after filtering the routine raises an
    error so that the caller can inspect upstream inputs.
  \item \textbf{Low-memory environments:} the per-worker temporary limit
    triggers the streaming fallback automatically; this may increase run time
    but prevents out-of-memory failures.
\end{itemize}

\section{Complexity and scaling}
Let $N_{\nu}$ be the number of photon-energy bins, $N_Z$ the number of charge
states retained, and $N_{\rm combos}$ the number of $(G_0,T,n_e)$ combinations.

\begin{itemize}
  \item Time complexity: building $R_{\rm pe}$ naively costs $O(N_Z N_{\nu})$;
    vectorized-chunked evaluation still costs the same arithmetic, but reduces
    memory overhead and can gain from optimized BLAS-like memory access.
  \item Memory complexity: the peak temporary memory is bounded by the chosen
    chunk size and by any cached kernels retained for reuse. The streaming
    fallback reduces peak memory to $O(N_{\nu,\text{chunk}})$.
\end{itemize}

\section{Practical recommendations}
\begin{itemize}
  \item On desktop/workstation with abundant RAM (e.g., >16~GB), increase the
    per-worker budget so the integrand can be built vectorized for speed.
  \item On HPC nodes with strict memory per rank, keep the per-worker budget
    small and rely on streaming to preserve stability.
  \item Persist computed per-(grain,a,rad_model) kernels between runs where
    possible (disk-backed cache) to avoid recomputation across experiments.
\end{itemize}

\section*{Acknowledgements and references}
The implementation follows the formalism of Weingartner & Draine (2001) for
photoelectric emission and grain charging, and adopts additional cross-section
and yield models described by Draine (1978, 2001). The interested reader
should consult Weingartner & Draine (2001) for a detailed derivation and
discussion of photoelectric yields, energy partitioning, and numerical
implementation details \cite{Weingartner2001}.

\begin{thebibliography}{9}
\bibitem{Weingartner2001}
J.~C. Weingartner and B.~T. Draine,
``Photoelectric Emission from Interstellar Dust: Grain Charging and Heating,''
\emph{Astrophysical Journal Supplement Series}, vol. 134, pp. 263--281, 2001.
% Additional references (Draine 1978, Draine 2001) may be added here as
% bibitems if full citation details are desired in the compiled document.
\end{thebibliography}

\end{document}
